% \iffalse meta-comment
%
% Copyright (c) 2011-2012 by Michael Helmling
%
% This file may be distributed and/or modified under the
% conditions of the LaTeX Project Public License, either
% version 1.3c of this license or (at your option) any later
% version. The latest version of this license is in:
% http://www.latex-project.org/lppl.txt
% and version 1.3c or later is part of all distributions of
% LaTeX version 2011/06/27 or later.
% \fi
%
% \iffalse
%<package>\NeedsTeXFormat{LaTeX2e}[2011/06/27]
%<package>\ProvidesPackage{agopt_ex}
%<package>  [2012/10/30 v0.4.1 optional points arg for solution environment]
%
%<*driver>
\documentclass{ltxdoc}
\setlength{\parindent}{0pt}
\usepackage{sectsty}
\allsectionsfont{\sffamily}
\usepackage[utf8]{inputenc}
\let\savetitle=\maketitle
\usepackage[english,classic,solution]{agopt_ex}
\let\maketitle=\savetitle
\providecommand*\pkg[1]{\textsf{#1}}
\providecommand*\env[1]{\texttt{#1}}
\providecommand*\mode[1]{\texttt{[#1]}}
\renewcommand{\nobreakspace}{\nobreak\ }
\Lecture{Using the \pkg{agopt\_ex} package}
\Semester{Winter 2012}
\Deadline{24.12.2039}
\IssueDate{10.02.2012}
\Lecturer{N.\,N.}
\Operator{Michael Helmling}
\renewcommand{\exercisesheettext}{Package Documentation}
\renewcommand{\withsolutiontext}{plus code}
\LectureShort{The \pkg{agopt\_ex} Package}

\cfoot{Documentation}

\EnableCrossrefs
\CodelineIndex
\RecordChanges
\begin{document}
\DocInput{agopt_ex.dtx}
\end{document}
%</driver>
% \fi

% \CheckSum{0}
% \changes{v0.4.1}{2012/10/30}{optional points arg for solution environment}
% \changes{v0.4}{2012/10/17}{add Koblenz mode}
% \changes{v0.3.1}{2012/02/10}{Fixed modern layout, added URL to documentation}
% \changes{v0.3}{2012/01/16}{A first complete proof-read, again lots of small changes}
% \changes{v0.2.2}{2012/01/11}{Improve on AG logo}
% \changes{v0.2.1}{2012/01/10}{Add font definitions}
% \changes{v0.2}{2012/01/06}{Largely rewritten}
% \changes{v0.1}{2011/12/22}{Initial version}
% \GetFileInfo{agopt_ex.sty}

% \title{The \textsf{agopt\_ex} package\thanks{This document
% corresponds to \textsf{agopt\_ex}~\fileversion,
% dated~\filedate. Obtain the newest version at \url{http://github.com/supermihi/latex}}}
% \author{ Michael Helmling \\ \texttt{helmling@mathematik.uni-kl.de}}
%

% \maketitle
% \begin{abstract}
%   The \pkg{agopt\_ex} package is an aid to generate exercise sheets for the Optimization Research Group, TU Kai\-sers\-lau\-tern, or optionally the Mathematisches Institut der Universität Koblenz.
% \end{abstract}
%
% The \pkg{agopt\_ex} package defines:
% \begin{itemize}
%   \item environments for exercises and solutions,
%   \item two layout variants (|classic| and |modern|, respectively,
%       the latter containig a colored logo); includes a nice footer and predefined macros for a ``in-class'' and ``take-home'' sections,
%   \item two locations (Kaiserslautern and Koblenz),
%   \item various ways to decide whether or not the solutions should be included in the output, including an \emph{auto-magic}$\textsuperscript{\textregistered}$ mechanism,
%   \item a number of useful features and macros.
% \end{itemize}
%
% This package documentation shows how to use the package by describing all package options and (re)defined macros. The corresponding source
% code snippets are included at the appropriate place for easy customization (and, of course, for instructional reasons).
%
% \section{Package Loading}
%
%
% \subsection{Package Options}\label{sec:options}
%
% \subsubsection*{Language Settings}
%  \begin{macro}{german}
%  \begin{macro}{english}
%  Define the language of the exercise sheet. The default is \texttt{german}. This option influences
%  various textual elements of the exercise sheet.
%    \begin{macrocode}
\def\ublanguage{german}
\DeclareOption{german}{
  \def\ublanguage{german}
}
\DeclareOption{english}{
  \def\ublanguage{english}
}
%    \end{macrocode}
%  \end{macro}
%  \end{macro}
%
% \subsubsection*{Toggle Solution Output}
%  \begin{macro}{solution}
%  \begin{macro}{nosolution}
%  These options define whether or not solutions should be included in the output document or not.
%  If neither option is present, \emph{auto-magic} detection is enabled.
%
%  With \emph{auto-magic} detection, the solutions are output if and only if the jobname contains the string
% ``olution'' (in |english| mode) or ``oesung'' (in |german| mode). Note that this is
% not a typo; the first letter is omitted in order to be case insensitive. If you need a different detection string,
% redefine the |\solutionfilename| macro.
% 
% The jobname is normally the name of the source file without the |.tex| extension, but can
% be overridden in (pdf)latex, xelatex etc.\ with the |-jobname=NAME| option. This allows
% for a convenient workflow: Create a single |.tex| file, e.\,g.\ |exercise1.tex|, then run
%
% \begin{itemize}
% \item \texttt{xelatex exercise1}
% \item \texttt{xelatex -jobname=solution1 exercise1}
% \end{itemize}
% (substitude \texttt{xelatex} with your favourite \LaTeX{} engine, e.\,g.\ \texttt{pdflatex}, \texttt{latex}, \dots)
% in order to generate \texttt{exercise1.pdf} (without solutions) and \texttt{solution1.pdf} (including solutions).
% The bash script |xeloetex| distributed together with this package shows how to merge both steps into one command
% that can be used as compile command in your favourite \TeX{} editor.
%    \begin{macrocode}
\newif\ifautoshowanswers
\newif\ifshowanswers
\showanswersfalse
\autoshowanswerstrue
\DeclareOption{solution}{\showanswerstrue\autoshowanswersfalse}
\DeclareOption{nosolution}{\showanswersfalse\autoshowanswersfalse}
%    \end{macrocode}
%  \end{macro}
%  \end{macro}
% \subsubsection*{Choose Location}
%  \begin{macro}{kaiserslautern}
%  \begin{macro}{koblenz}
%  Selects the location, where Kaiserslautern is the default.
%    \begin{macrocode}
\newif\ifkoblenz
\koblenzfalse
\DeclareOption{koblenz}{\koblenztrue}
\DeclareOption{kaiserslautern}{\koblenzfalse}
%    \end{macrocode}
%  \end{macro}
%  \end{macro}
%
% \subsubsection*{Choose Layout}
%  \begin{macro}{classic}
%  \begin{macro}{modern}
%  Defines the style of the exercise sheet. \texttt{modern} uses a colored graphical logo of the AG in the title
% (as in this document). \texttt{classic} resembles the classical exercise sheet style which hasn't changed for the
% past 30 years.
%    \begin{macrocode}
\newif\ifmodern
\DeclareOption{modern}{\moderntrue}
\DeclareOption{classic}{\modernfalse}
\moderntrue % the default
%    \end{macrocode}
%  \end{macro}
%  \end{macro}
% This closes the options section.
%    \begin{macrocode}
\ProcessOptions\relax
%    \end{macrocode}
% \subsection{Fonts}
% The package configures \TeX{} to use fonts of the Linux Libertine \ifxetex\char"E040\else\libertineGlyph{uniE040}\fi{} family and the Euler math font. The implementation
% differs for |(pdf)latex| and |xelatex|. For this package to work with |xelatex|, you need to have the Linux Libertine
% and Linux Biolinum OpenType fonts installed.
%    \begin{macrocode}
\RequirePackage{ifxetex}
\RequirePackage{ifthen}
\ifxetex
  \RequirePackage{amsfonts,amssymb}
  \RequirePackage{euler}
  \RequirePackage{xltxtra}
  \RequirePackage{xunicode}
  \RequirePackage{polyglossia}
  \defaultfontfeatures{Mapping=tex-text} % needed for -- and --- to work
  \setromanfont[Numbers=Proportional]{Linux Libertine O}
  \setsansfont[Numbers=Proportional]{Linux Biolinum O}
  \ifthenelse{\equal{\ublanguage}{german}}{
	  \setdefaultlanguage{german}
  }{
	   \setdefaultlanguage[variant=american]{english}
	}
\else
  \ifthenelse{\equal{\ublanguage}{german}}{
  	\usepackage[ngerman]{babel}
  }{
  	\usepackage[american]{babel}
  }
  \RequirePackage{libertine}
  \RequirePackage[T1]{fontenc}

  \usepackage{euler}
\fi
%    \end{macrocode} 
% \subsection{Required Packages}
% The following packages are needed by \pkg{agopt\_ex}:
%    \begin{macrocode}
\RequirePackage{amsmath}
\RequirePackage{geometry}
\RequirePackage{hyperref}
\RequirePackage{fancyhdr}
\RequirePackage{zref-totpages}
\RequirePackage{prettyref}
\RequirePackage{url}
\RequirePackage{xspace}
%    \end{macrocode}
%
% \section{Providing Lecture and Exercise Parameters}
%
%
% The following lecture and tutorial data should be set in every exercise sheet.
%  \begin{macro}{\Lecture}
%   Specify the name of the lecture (e.\,g. ``Praktische Mathematik: Lineare und Netzwerkoptimierung'').
%    \begin{macrocode}
 \def\Lecture#1{\def\lecture{#1}}
%    \end{macrocode}
%  \end{macro}
%  \begin{macro}{\LectureShort}
%   Specify a short name of the lecture, used in the footer (e.\,g. ``PraMa Optimierung'').
% \begin{macrocode}
\def\LectureShort#1{\def\lectureshort{#1}}
%    \end{macrocode}
%  \end{macro}
%  \begin{macro}{\Sheetnumber}
%   Specify the exercise sheet number.
% \begin{macrocode}
\def\Sheetnumber#1{\def\sheetnumber{#1}}
%    \end{macrocode}
%  \end{macro}
%  \begin{macro}{\Deadline}
%   Specify the deadline for turn-in exercises. May include additional information such as ``in the lecture'' or ``into the mailboxes in building 48''.
% \begin{macrocode}
\def\Deadline#1{\def\deadline{#1}} 
%    \end{macrocode}
%  \end{macro}
%  \begin{macro}{\IssueDate}
%   Specify the date when the sheet was issued.
% \begin{macrocode}
\def\IssueDate#1{\def\issuedate{#1}} 
%    \end{macrocode}
%  \end{macro}
%  \begin{macro}{\Lecturer}
%   Specify the name of the lecturer.
% \begin{macrocode}
\def\Lecturer#1{\def\lecturer{#1}} 
%    \end{macrocode}
%  \end{macro}
%  \begin{macro}{\Operator}
%   Specify the name of the exercise operator.
% \begin{macrocode}
\def\Operator#1{\def\operator{#1}} 
%    \end{macrocode}
%  \end{macro}
%  \begin{macro}{\Semester}
%   Specify the current semester or term (e.\,g. ``winter term 2012'').
% \begin{macrocode}
\def\Semester#1{\def\semester{#1}}
%    \end{macrocode}
%  \end{macro}
%  \begin{macro}{\Homepage}
%  This optional parameter defines a homepage for the exercises. If it is
%  used, the document output will contain a note where to download exercises.
%    \begin{macrocode}
\def\Homepage#1{\def\homepage{#1}}
%    \end{macrocode}
%  \end{macro}
%  \begin{macro}{\InclassDate}
%  This optional parameter defines the date for in-class exercises.
%    \begin{macrocode}
\def\InclassDate#1{\def\inclassdate{#1}}
%    \end{macrocode}
%  \end{macro}
% The parameters defined by the above macros can be accessed by their lowercase equivalents.
%    \begin{macrocode}
\def\lecture{Default lecture name}
\def\lectureshort{PraMa Optimierung}
\def\sheetnumber{1}
\def\deadline{}
\def\issuedate{06.12.1970}
\def\lecturer{Lecturer}
\def\operator{Exercise Operator}
\def\semester{Semester}
\def\homepage{}
%    \end{macrocode}
%
%
% \subsection{Change Default Textual Elements}
%
% The words used for ``Exercise'', ``Sheet'' etc.\ can be modified
% by redefining the following commands:
%    \begin{macrocode}
\ifthenelse{\equal{\ublanguage}{german}}{
  \def\solutiontext{L\"osung}
  \def\exercisetext{Aufgabe}
  \newcommand{\exercisesheettext}{\"Ubungsblatt}
  \def\withsolutiontext{mit L\"osung}
  \def\pagetext{Seite}
  \def\pointstext{Punkte}
  \def\solutionsheettext{L\"osungsblatt}
  \def\deadlinetext{Abgabe bis}
  \def\solutionfilename{oesung}
  \def\lecturetext{Vorlesung}
  \newcommand{\exercisestext}{\"Ubungen}
  \newcommand{\homepagetext}{Dieses \"Ubungsblatt sowie weitere %
  Informationen zur \"Ubung sind unter \url{\homepage} erh\"altlich.}
  \newcommand{\inclasstexttitle}{Pr\"asenz\"ubungen}
  \newcommand{\inclasstext}{Zur Bearbeitung in der \"Ubung am \inclassdate}
  \newcommand{\takehometexttitle}{Haus\"ubungen}
  \newcommand{\takehometext}{Bitte bis \deadline{} abgeben.}
  \newcommand{\deadlinepre}{\textbf{Abgabefrist: }}
}{
  \def\solutiontext{Solution}
  \def\exercisetext{Exercise}
  \def\exercisesheettext{Exercise Sheet}
  \def\solutionsheettext{Solution Sheet}
  \def\withsolutiontext{including solutions}
  \def\pagetext{Page}
  \def\pointstext{points}
  \def\deadlinetext{Due date:}
  \def\solutionfilename{olution}
  \def\lecturetext{Lecture}
  \def\exercisestext{Exercises}
  \newcommand{\homepagetext}{Download of exercises at \url{\homepage}}
  \newcommand{\inclasstexttitle}{In-Class Exercises}
  \newcommand{\inclasstext}{To be done in the tutorial on \inclassdate}
  \newcommand{\takehometexttitle}{Turn-In Exercises}
  \newcommand{\takehometext}{Please hand in by \deadline{}}
  \newcommand{\deadlinepre}{\textbf{Deadline: }}
}
%    \end{macrocode}
% For example, if you wish to name exercises ``Problem'' rather than ``Exercise'',
% simply put \[|\renewcommand{\exercisetext}{Problem}|\] in your preamble.
%
%
% \section{Typesetting Exercises and Solutions}
%
% \subsection{Exercises}
%
% \DescribeEnv{exercise} The |exercise| environment is used in the
% following way:
% 
% \noindent|\begin{exercise}| \oarg{title} \marg{points}\\
% |  ...|\\
% |\end{exercise}|
% 
% The parameter \meta{points} will be typeset in
% parenthesis after the exercise title, unless it is empty.
% If the optional \meta{title} is given, the exercise title is typeset after the
% exercise number, separated by an endash (--).
%
% Exercises are numbered by a special counter (|exercise|); the number is
% displayed in the style |x.y| where |x| is the sheet number and
% |y| the exercise number on the sheet. You can thus use |\label| and |\ref|
% for exercise refercening as well as |\theexercise| to output the current
% exercise number.
%
% As an example, the code
% \begin{verbatim}
% \begin{exercise}[$P \neq NP$]{4}
%   Prove that $P$ is a proper subset of $NP$.
% \end{exercise}
% \end{verbatim}
%
% will be output as
%
% \begin{exercise}[$P \neq NP$]{4}
%   Prove that $P$ is a proper subset of $NP$.
% \end{exercise}
%
%
%    \begin{macrocode}
\newcommand{\exheader}[1]{\par\vspace{2.5mm}\noindent{\bfseries #1}\par\vspace{1.5mm}}
\newcounter{exercise}
\setcounter{exercise}{0}
\newenvironment{exercise}[2][{}]%
{%
  \refstepcounter{exercise}
  \exheader{\exercisetext{} \sheetnumber.\arabic{exercise}
  \ifthenelse{\equal{#1}{}}{}{-- #1}
  \ifthenelse{\equal{#2}{}}{}{(#2 \pointstext)}}
}%
{\par\vspace{2mm}}
%    \end{macrocode}
%
% Subexercises can be typeset with usual |\enumerate| environments. In order not
% to mix up exercise and subexercise numbering, this package sets the first-order
% enumeration labelling to alphabetic numbering and the second order to arabic:
%    \begin{macrocode}
\RequirePackage{enumitem}
\setlist[enumerate,1]{label=\alph*)}
\setlist[enumerate,2]{label=\arabic*.}
%    \end{macrocode}
%
% \subsection{Solutions}
%
% \DescribeEnv{solution} The |solution| environment can be used to create a
% sample solution. You can decide whether or not solutions will be included
% in the output, in order to distinguish between exercise and solution sheets
% (see Section~\ref{sec:options}).
%
% The |solution| environment is used as follows:
%
% \noindent|\begin{solution}| \oarg{points}\\
% |  ...|\\
% |\end{solution}|
%
% The optional \meta{points} parameter is typeset in the same way as the \meta{points} argument of the \meta{solution} environment. It may be used to denote the point split in case of subexercises.
%
% For example, the code
% \begin{verbatim}
% \begin{solution}[2+2]
%   Base clause: Let $N=1$, then obviously $P=NP$. 
% \end{solution}
% \end{verbatim}
%
% will be output to (if solution output is active)
%
% \begin{solution}[2+2]
%   Base clause: Let $N=1$, then obviously $P=NP$. 
% \end{solution}
%
%
% \subsection{Implementation of the Auto-Magic Solution Feature}
% If neither |solution| nor |nosolution| is provided as package option,
% test if the |\jobname| contains the (language specific) word
% for ``solution''. The test requires the \pkg{xstring} package.
%    \begin{macrocode}

\ifautoshowanswers
\RequirePackage{xstring}
\IfSubStr*{\jobname}{\solutionfilename}{
 \showanswerstrue
}{
 \showanswersfalse
}
\fi
\newenvironment{solution}[1][{}]%
{%
  \ifshowanswers
    \exheader{\solutiontext{} \sheetnumber.\arabic{exercise}%
    \ifthenelse{\equal{#1}{}}{}{ (#1 \pointstext)}:}
  \else
    \par\vspace*{0pt}%
    \setbox\z@\vbox\bgroup
  \fi
}{%
  \ifshowanswers
    %
  \else
    \egroup
  \fi
}%
%    \end{macrocode}
%
%
% \subsection{In-Class and Take-Home Exercises}
%
%\begin{macro}{\inclass}
%\begin{macro}{\takehome}
% These optional macros create a title that marks the begin of the ``in-class'' or ``take-home'' part, respectively,
% of the exersice sheet.
%    \begin{macrocode}
% marks if an exercise type (inclass, takehome) was explicitly chosen, because otherwise the layout
% has to automatically print the deadline information.
\newif\ifexplicittype 
\explicittypefalse
\newcommand{\inclass}{\par{\large
\ifmodern\textsc{\inclasstexttitle}\else\MakeUppercase{\inclasstexttitle}\fi}\\
(\inclasstext)\par
\explicittypetrue
}
\newcommand{\takehome}{\par{\large
\ifmodern\textsc{\takehometexttitle}\else\MakeUppercase{\takehometexttitle}\fi}\\
(\takehometext)\par
\explicittypetrue
}
%    \end{macrocode}
%\end{macro}
%\end{macro}
% \section{Miscellaneous Features}
%
% \subsection{PDF parameters}
% This package sets some PDF parameters according to the exercise sheet
% definition.
%    \begin{macrocode}
\hypersetup{%
  pdftitle={\lecture, \exercisesheettext{} \sheetnumber}, %
  pdfauthor={\ifkoblenz Mathematisches Institut, Universität Koblenz\else Optimization Research Group, TU Kaiserslautern\fi}, %
  pdfcreator={\ifxetex XeLaTeX \else LaTeX2e \fi}}
%    \end{macrocode}
%
%
% \subsection{Referencing Exercises and Solutions}
%
% This package defines to reference formats for the \pkg{prettyref} package
% which can be used to reference exercises and solutions, respectively.
% Example:\\
% |Use the graph of \prettyref{ex:dijkstra} and ...|\\
% Would be typeset as, say,
% \begin{quote}Use the graph of~Exercise 2 and \dots\end{quote}
%
%
%    \begin{macrocode}
\newrefformat{ex}{\exercisetext~\ref{#1}}
\newrefformat{solution}{\solutiontext~\ref{#1}}
%    \end{macrocode}
%
% \subsection{Headers and Footers}
% \pkg{agopt\_ex} uses \pkg{fancyhdr} to set an empty header and a nice footer. 
% You can modify the following default layout if you wish.
%    \begin{macrocode}
\pagestyle{fancy}
\fancyhead{}
\renewcommand{\headrulewidth}{0pt}
\renewcommand{\footrulewidth}{.4pt}
\cfoot{\ifshowanswers\solutionsheettext{}\else\exercisesheettext{}\fi{} \sheetnumber}
\rfoot{\pagetext{} \thepage/\ztotpages}
\lfoot{\lectureshort}
%    \end{macrocode}
%
% \section{Implementation of the Layouts}
%
% The modern layout uses \pkg{tikz} to draw the logo.
%    \begin{macrocode}
\newcommand{\titledateline}{%
    \ifthenelse{\equal{\deadline}{}}%
        {\inclasstext}%
        {\deadlinetext{} \deadline{}}%
}
\ifmodern
	\RequirePackage{tikz}
	\definecolor{tublau}{rgb}{0.125,0.34,0.68}
	\renewcommand{\maketitle}{
	\hrule\vspace{2mm}
	\ifkoblenz
	    \begin{minipage}{0.65\textwidth}
	\else
    	\begin{minipage}{0.55\textwidth}
    \fi
	{\sffamily \lecture{}\;\textbullet\;\semester\\
	\LARGE \scshape  \exercisesheettext{} \sheetnumber %
	\ifshowanswers%
	  {\Large{} (\withsolutiontext)}%
	\fi\\
	\small \upshape \itshape \rmfamily \titledateline}
	\end{minipage}
	\ifkoblenz
    	\begin{minipage}{0.34\textwidth}
    \else
        \begin{minipage}{0.44\textwidth}
    \fi
	\begin{flushright}
	\ifkoblenz
	    \includegraphics[width=0.8\textwidth]{koblenzlogo}
	\else
		\begin{tikzpicture}[klumpen/.style={minimum size=4mm,rectangle},
		                    every edge/.append style={very thick},scale=.9]
		 \node[fill=red,klumpen] (k1) at (0,0) {};
		 \node[fill=tublau,klumpen] (k2) at (2,0) {} edge (k1);
		 \node[fill=tublau,klumpen] (k3) at (2,-1) {} edge(k2);
		 \node[fill=red,klumpen] (k4) at (5,-1) {} edge(k3);
		 \node[font={\sffamily\bfseries\fontsize{15}{16}\selectfont}] at (1,-.5) {OPT};
		 \node[font={\sffamily\fontsize{8}{7}\selectfont},anchor=west] at (2.3,-.5)
		   {\begin{minipage}{2.6cm}Optimization\\Research Group\end{minipage}};
		\end{tikzpicture}
	\fi
	\end{flushright}
	\end{minipage}\vspace{2mm}\hrule
	\begin{center}\small
	 \textbf{\lecturetext:} \lecturer\\
	 \textbf{\exercisestext:} \operator
	\end{center}\vspace{-2mm}
	\ifthenelse{\equal{\homepage}{}}{}{
	{\small \homepagetext}
	}
	}
%    \end{macrocode}
%
% This is the implementation of the classic layout.
%    \begin{macrocode}
\else 
	\renewcommand{\maketitle}{
	\begin{minipage}{0.49\textwidth}
	\begin{flushleft}
	\ifkoblenz
	    Universit\"at Koblenz-Landau, Campus Koblenz\\
	    Mathematisches Institut\\
	\else
    	Technische Universit\"at Kaiserslautern\\
    	Fachbereich Mathematik\\
    \fi
	\issuedate
	\end{flushleft}
	\end{minipage}
	\begin{minipage}{0.49\textwidth}
	\begin{flushright}
	\lecturer\\
	\operator\\
	\semester
	\end{flushright}
	\end{minipage}
	
	\begin{center}
	{\Large \bfseries \lecture}\\[0.6cm]
	{\Large \bfseries%
	\ifshowanswers%
	  \solutionsheettext{}%
	\else
	  \exercisesheettext{}%
	\fi{} \sheetnumber}\\[1cm]
	\end{center}
	}
	% URL at end of document
	\AtEndDocument{%
	\ifexplicittype
	\else
		\par
	  \deadlinepre\takehometext
	\fi
	\ifthenelse{\equal{\homepage}{}}{}{
		\begin{center}
			\vfill{\small \homepagetext}
		\end{center}
		}
	}
\fi
%    \end{macrocode}
% \PrintChanges
% \PrintIndex
% \Finale
